\subsection{Ising forests at every fixed finite temperature}
\label{sec:threshold-positive-forests}

\begin{theorem}[Forest simulation]
\label{thm:threshold-ising-forests}
Let \(G\) be any finite forest and
\[
 \mu(x)\propto
 \exp\left(\beta\sum_{uv\in E(G)}J_{uv}x_ux_v
                 +\beta\sum_v h_vx_v\right),\qquad
 |J_{uv}|\le J_{\mathrm{edge}}.
\]
The fields, signs of the couplings, degrees, component sizes, and
geometry may be arbitrary. Fix \(\beta J_{\mathrm{edge}}<\infty\). Then this family
has full AA with constants depending only on \(\beta J_{\mathrm{edge}}\), uniformly
over the fields and forests. The conclusion also holds uniformly
after pinning any set of spins.

More explicitly, put \(\rho=1-e^{-4\beta J_{\mathrm{edge}}}\). For \(\rho>0\),
choose
\[
 L=\max\left\{1,
 \left\lceil\frac{2\log(8T/e)}{\log(1/\rho)}\right\rceil\right\}.
\]
For \(\rho=0\), take \(L=1\). A valid query bound is
\[
 2S_B\bigl((4L+2)T,e/2;2\bigr).
\]
In particular, there is no dependence on \(n\) or \(\log n\).
\end{theorem}

\begin{proof}
Root each component. A positive distribution that is Markov on a tree
has the directed factorization
\[
 \mu(x)=\prod_{r\text{ root}}\pi_r(x_r)
          \prod_{e=(u,v)}K_e(x_u,x_v).
\]
For this Ising law, summing over the descendants of \(v\) gives
\[
 K_e(s,t)=\frac{\exp(t[a_v+\beta J_es])}
                    {2\cosh(a_v+\beta J_es)}
\]
with an effective real field \(a_v\). It may be large and need not be
uniformly bounded.

For each column \(t\in\{-1,+1\}\), define
\[
 m_e(t)=\min_{s\in\{-1,+1\}}K_e(s,t),\qquad
 q_e=\sum_t m_e(t),\qquad \lambda_e=1-q_e.
\]
Finite fields and finite couplings imply \(q_e>0\), so the reset law
below is defined. If \(\lambda_e>0\), the two rows of
\((K_e(s,t)-m_e(t))/\lambda_e\) are point masses on opposite spins.
Thus there is a known \(\sigma_e\in\{-1,+1\}\) such that the kernel
can be generated by independent source bits
\[
 C_e\sim\operatorname{Bernoulli}(\lambda_e),\qquad
 Y_e\sim m_e/q_e,
\]
and the rule
\[
 X_v=\begin{cases}
       \sigma_eX_u,& C_e=1,\\
       Y_e,& C_e=0.
      \end{cases}
\]
When \(\lambda_e=0\), use only the second branch. Independently
generate the root spins with laws \(\pi_r\). All edge pairs
\((C_e,Y_e)\) and root bits are mutually independent. This is a
product of blocks of size at most two, so
Theorem~\ref{thm:threshold-block-products} supplies its full-AA
simulator with constants independent of every source bias.

The posterior of these sources has a special locality property.
Conditioned on the complete observed configuration \(X=x\), the root
source bits equal \(x_r\), and the edge-source pairs remain independent
across edges. The pair on \(e=(u,v)\) depends only on \((x_u,x_v)\).
If \(x_v\ne\sigma_ex_u\), then \(C_e=0\) and \(Y_e=x_v\).
Otherwise,
\[
 \Pr(C_e=1\mid X=x)
   =\frac{\lambda_e}{K_e(x_u,x_v)}
   =1-\frac{\min_sK_e(s,x_v)}{\max_sK_e(s,x_v)}.
\]
On the copy branch \(Y_e\) retains its original law; on the reset
branch it is fixed to \(x_v\). These formulas follow by factoring the
prior source law together with the local deterministic constraints
defining \(X\). In particular, they do not assume that conditioning
preserves the original source independence without checking it.

The Ising kernel has, for either column,
\[
 \frac{\min_sK_e(s,t)}{\max_sK_e(s,t)}
       \ge e^{-4\beta|J_e|}\ge e^{-4\beta J_{\mathrm{edge}}}.
\]
For example, when \(J_e\ge0\),
\(K_e(+,+)/K_e(-,+)
 =e^{2\beta J_e}\cosh(a_v-\beta J_e)/\cosh(a_v+\beta J_e)
 \le e^{4\beta J_e}\); the other column and signs follow in the same
way. Hence \(\Pr(C_e=1\mid X=x)\le\rho<1\), uniformly in \(x\),
the fields, and the graph.

Define a truncated source map \(F_L\) for each vertex by examining
its ancestor path for at most \(L\) edges. If a reset is encountered,
its spin together with the intervening known signs determines the
output. If the component root is reached first, use its source spin.
If neither happens within \(L\) edges, output an arbitrary fixed spin.
This map can be evaluated using at most \(2L+1\) source bits; a
coherent query costs at most \(4L+2\) source queries.

Given \(X=x\), failure of \(F_L\) at a vertex of depth at least
\(L\) is possible only if all \(L\) inspected edges use their copy
branch. Posterior edge independence yields
\[
 \Pr[F_L(Z)_v\ne x_v\mid X=x]\le\rho^L.
\]
For shallower vertices the map is exact. To supply any source query
from its posterior, read the two endpoints of the relevant edge
(or the root spin), generate the conditional source pair, and cache
it. Consequently \(s\) source-bit requests need at most \(2s\)
physical-spin queries.

Apply Proposition~\ref{prop:threshold-local-posterior} with source
error \(e/2\). The total error is at most
\(e/2+4T\rho^{L/2}\le e\), and the displayed query count follows.
Finally, pinning vertices deletes those vertices from the unqueried
forest and adds fields \(J_{uv}x_u\) to their unpinned neighbors.
The remaining couplings are unchanged, and the proof is uniform over
all fields. This proves the pinning assertion.
\end{proof}

This theorem concerns the classical diagonal Gibbs law. It does not
automatically apply to the measurement distribution of a noncommuting
transverse-field Hamiltonian on the same tree: imaginary-time
trajectories generally destroy this simple binary Markov
factorization. Also, the displayed constants deteriorate as
\(\beta J_{\mathrm{edge}}\to\infty\); the theorem alone does not establish a
uniform result for arbitrary dimension-dependent inverse temperatures.
